\chapter*{Appendix}
\addcontentsline{toc}{chapter}{Appendix}
\setcounter{table}{0}
\renewcommand{\thetable}{A.\arabic{table}}

\section*{A. Judge prompt}\label{app:prompt}
\addcontentsline{toc}{section}{A. Judge prompt}

The user turn of the judge contains the image followed by the text below, where \texttt{\{question\}} and \texttt{\{answer\}} are replaced by the record's question and answer. The assistant turn is the JSON array described in Section~\ref{sec:judge}.

\begin{small}
\begin{verbatim}
Identify all hallucinated text spans in the vision model's response.
A hallucination is text that is factually incorrect compared to what
is actually visible in the image.

Categories:
- invention: entities or facts completely absent from the image
- mischaracterization: visible content described incorrectly
- OCR: text visible in the image is misread or mistranscribed
- miscounting: wrong quantity of visible items
- other: hallucination not fitting the above

List spans in the order they appear in the response, left to right.

Return ONLY a JSON array, no markdown fences, no extra text:
[{"text": "<exact text copied verbatim from the response>",
  "prob": <0.0-1.0>, "label": "<category>"}]
If there are no hallucinations, return: []

Question asked: {question}

VLM response to analyse:
{answer}
\end{verbatim}
\end{small}

\section*{B. Hyperparameters}
\addcontentsline{toc}{section}{B. Hyperparameters}

\begin{center}
    \captionof{table}{Fixed hyperparameters. None was tuned on validation; $k$ was chosen once among three values on val\_select.}
    \label{tab:hparams}
    \resizebox{\textwidth}{!}{%
    \begin{tabular}{lll}
    \toprule
    Component & Parameter & Value \\
    \midrule
    Judge (QLoRA) & quantisation & 4-bit NF4, frozen backbone \\
     & LoRA & rank 16, $\alpha = 16$, language-model attention and MLP projections \\
     & optimiser & AdamW, lr $2 \cdot 10^{-4}$, weight decay 0.01, grad.\ clip 1.0 \\
     & schedule & 50 warm-up steps, linear decay \\
     & batch & 1 $\times$ 8 accumulation steps, 2 epochs (3,020 steps) \\
     & max.\ sequence & 4,096 tokens \\
    Decoding & greedy & 1 decode \\
     & sampling & 8 decodes, temperature 0.7 \\
    Vote & threshold & $k = 3$ (of 9 or 10 voters), per label \\
     & output probability & per character, rounded to 2 decimals \\
    Probe features & frozen model & Qwen3-VL-8B-Instruct (bf16), teacher-forced \\
     & per token & entropy, NLL, look-back ratio ($36 \times 32$), last hidden state (4,096) \\
     & reduction & PCA on train: hidden $\to$ 256, look-back $\to$ 128; standardised \\
    Probe model & encoder & Linear 512 $\to$ BiGRU 256, language embedding 8, dropout 0.2 \\
     & heads & 5 binary linear-chain CRFs, multi-label targets \\
     & training & Adam, lr $10^{-3}$, batch 16, $\le$ 20 epochs, patience 3 on val\_select loss \\
     & max.\ length & 1,024 tokens \\
    Seeds & all trained components & 42, 123, 2024 \\
    \bottomrule
    \end{tabular}}
\end{center}

\section*{C. Per-cell results for H4}
\addcontentsline{toc}{section}{C. Per-cell results for H4}

\begin{center}
    \captionof{table}{Per-category Cor on val\_select for each language (three-seed means): self-consistency voter F2 / probe alone. Spearman correlation over the 20 cells: 0.95.}
    \label{tab:h4cells}
    \begin{tabular}{lcccc}
    \toprule
    Category & en & fr & it & zh \\
    \midrule
    Invention & 0.172 / 0.119 & 0.215 / 0.186 & 0.247 / 0.158 & 0.290 / 0.216 \\
    Mischaracterization & 0.149 / 0.070 & 0.183 / 0.102 & 0.212 / 0.110 & 0.144 / 0.038 \\
    OCR & 0.448 / 0.429 & 0.360 / 0.283 & 0.304 / 0.158 & 0.225 / 0.166 \\
    Miscounting & 0.466 / 0.362 & 0.531 / 0.475 & 0.446 / 0.359 & 0.376 / 0.367 \\
    Other & 0.130 / 0.050 & 0.022 / 0.032 & 0.076 / 0.090 & 0.008 / 0.041 \\
    \bottomrule
    \end{tabular}
\end{center}

\section*{D. Reproducibility}
\addcontentsline{toc}{section}{D. Reproducibility}

All code, the frozen validation split, the decision log, the pre-registered rules (R1--R6) and every validation score file are kept in the project repository under \texttt{thesis\_v2/}. The main entry points are:
\begin{itemize}
    \item \texttt{part1\_models/train.py}, \texttt{predict.py}: judge training and decoding;
    \item \texttt{common/span\_io.py}: the label-aware span decoder;
    \item \texttt{common/voting.py}: the per-label vote, with graded output;
    \item \texttt{part3\_uncert/crf\_probe.py}: the probe;
    \item \texttt{integration/build\_candidates.py}, \texttt{graded.py}: fusion and the final decisions;
    \item \texttt{common/evaluate.py}: the official scoring with the coverage guard and the paired bootstrap.
\end{itemize}
Official test results are recorded in \texttt{results/test\_results.txt}.
% TODO: add the public repository URL if the code is released.
